\clearpage
\subsection{再论 Koszul 对偶集}\label{sec:5o}
设 $\A^{\to}$ 是有向 $\Ai$ 范畴，$\A=\Tw\A^{\to}$。回顾\stref{引理}{3.36}，存在典范的上同调全忠实 $\Ai$ 函子 $\widetilde{\mathcal I}:\A\to\Q=\modcat(\A^{\to})$。由定义，$Y^!_{m+1-k}\in\Ob\A$ 在此函子下的像是 $\Ai$ 模 $\Y^!_{m+1-k}=\mathcal T_{Y_1}\cdots\mathcal T_{Y_{k-1}}(\Y_k)$。
\SourcePage{93}
这样的模的具体形式很容易写出。基础分次向量空间为
\begin{equation}\label{eq:5.35}
\begin{aligned}
\Y^!_{m+1-k}(Y_l)=\bigoplus_r\bigoplus_{i_1,\ldots,i_r}
\bigl(&\hom_{\A}(Y_{i_r},Y_k)\otimes\hom_{\A}(Y_{i_{r-1}},Y_{i_r})\otimes\cdots\\[-2pt]
&\otimes\hom_{\A}(Y_{i_1},Y_{i_2})\otimes\hom_{\A}(Y_l,Y_{i_1})\bigr)[r],
\end{aligned}
\end{equation}
其中对 $r\geq0$、$l\leq i_1<i_2<\cdots<i_r<k$ 求和；特别地，$r=0$ 的直和项就是 $\hom_{\A}(Y_l,Y_k)$。其次，
\begin{equation}\label{eq:5.36}
\begin{aligned}
&\mu^1_{\Y^!_{m+1-k}}(b_{r+1}\otimes\cdots\otimes b_1)\\
&\quad=\sum_{p,n}(-1)^{\dagger}b_{r+1}\otimes\cdots\otimes b_{n+p+1}
\otimes\mu^p_{\A}(b_{n+p},\ldots,b_{n+1})\otimes b_n\otimes\cdots\otimes b_1,
\end{aligned}
\end{equation}
其中 $\dagger=|b_1|+\cdots+|b_n|-n$；对于 $d\geq2$，
\[
\begin{aligned}
&\mu^d_{\Y^!_{m+1-k}}(b_{r+1}\otimes\cdots\otimes b_1,a_{d-1},\ldots,a_1)\\
&\quad=\sum_{n\geq1}b_{r+1}\otimes\cdots\otimes b_{n+1}\otimes\mu^{d-1+n}_{\A}(b_n,\ldots,b_1,a_{d-1},\ldots,a_1).
\end{aligned}
\]
\begin{definition*}
现在考虑\emph{简单 $\Ai$ 模} $\mathit{SS}_k$：令 $\mathit{SS}_k(Y_k)=\K$，对所有其他 $l$ 令 $\mathit{SS}_k(Y_l)=0$，并取 $\mu^2_{\mathit{SS}_k}(\cdot,e_{Y_k})=\Id_{\K}$，其余结构映射均为零。
\end{definition*}
存在典范模同态
\[
t:\Y^!_{m+1-k}\longrightarrow\mathit{SS}_k.
\]
一阶部分 $t^1$ 为从 $\Y^!_{m+1-k}(Y_k)=\hom_{\A}(Y_k,Y_k)=\K\cdot e_{Y_k}$ 到 $\mathit{SS}_k(Y_k)$ 的恒等映射；高阶映射 $t^d$（$d\geq2$）恒为零。\znote{85}{简单模 $\mathit{SS}_k$ 是有向 $\Ai$ 范畴 $\A^{\to}$ 的模范畴中最小的对象：它只在 $Y_k$ 处取值 $\K$，其余为零，且 $\mu^2(\cdot,e_{Y_k})$ 为恒等。引理 5.24 说 Koszul 对偶对象 $Y^!_{m+1-k}$ 的 Yoneda 模同构于 $\mathit{SS}_k$，因此 Koszul 对偶集恰好给出一族简单模；这也是后文把 $\Q$ 与“每个对象都由简单模经迭代锥得到”这一结构联系起来的出发点。[来源：并见 Seidel, Fukaya categories and Picard–Lefschetz theory, §5o（引理 5.24）与 §5p。]}
\begin{lemma}\label{st:5.24}
$[t]$ 是 $H^0(\Q)$ 中的同构。
\end{lemma}
\begin{proof}
复形\eref{5.35}具有显然的递增长度滤过 $F^r$。将 $r$ 换成 $-r$，使其变为递降滤过，考察相应的谱序列，它收敛到 $H(\Y^!_{m+1-k}(Y_l))$。起始项为
\[
E_1^{-r,s}=\left(\bigoplus_{i_1,\ldots,i_r}\Hom_{H(\A)}(Y_{i_r},Y_k)\otimes\cdots\otimes\Hom_{H(\A)}(Y_l,Y_{i_1})\right)^s,
\]
其中 $i_k$ 如\eref{5.35}，微分为
\[
\partial^{-r,s}(b_{r+1}\otimes\cdots\otimes b_1)=\sum_n(-1)^{\dagger+|b_{n+1}|}b_{r+1}\otimes\cdots\otimes b_{n+2}b_{n+1}\otimes\cdots\otimes b_1,
\]
\SourcePage{94}
其中 $\dagger$ 如\eref{5.36}。这正是一个约化条形分解；更准确地说，若对所有 $l$ 取直和，所得就是分次 $H(\A)$ 模的 Abel 范畴 $\operatorname{Mod}^0(H(\A))$ 中，简单对象 $H(\mathit{SS}_k)$ 的标准投射分解。因此，当 $k=l$ 时，谱序列的 $E_2$ 项为一维，由 $[e_{Y_k}]$ 张成；其余情形为零。与 $\mathit{SS}_k(Y_l)$ 的平凡滤过作比较很直接，结论是对任意 $l$，$t^1:\Y^!_{m+1-k}(Y_l)\to\mathit{SS}_k(Y_l)$ 都在上同调上诱导同构。再由\stref{引理}{1.16}，得到所需结果。
\end{proof}
\begin{note}\label{st:5.25}
概述一下这个观点如何给出谱序列\eref{5.28}的另一种构造。任意 $\mathcal U\in\Ob\Q$ 都带有自然的子模递降滤过
\[
(F^r\mathcal U)(Y_k)=\begin{cases}\mathcal U(Y_k),&k\leq m-r,\\0,&\text{其余情形}.\end{cases}
\]
分次部分为 $F^r\mathcal U/F^{r+1}\mathcal U\cong\mathcal U(Y_{m-r})\otimes\mathit{SS}_{m-r}$，其中将 $\mathcal U(Y_{m-r})$ 视为向量空间的链复形。现在取 $\mathcal U=\mathcal U_1$，考虑 $\hom_{\Q}(\mathcal U_0,\mathcal U_1)$ 上的诱导滤过。所得谱序列的起始项为
\begin{equation}\label{eq:5.37}
E_1^{rs}=\bigl(H(\hom_{\Q}(\mathcal U_0,\mathit{SS}_{m-r}))\otimes H(\mathcal U_1(Y_{m-r}))\bigr)^{r+s}.
\end{equation}
现在假设 $\mathcal U_0,\mathcal U_1$ 是扭曲复形在上同调全忠实嵌入 $\widetilde{\mathcal I}:\A\to\Q$ 下的像。利用\stref{引理}{2.12}与\stref{引理}{5.24}，就可以将\eref{5.37}化为\eref{5.28}的形式。
\end{note}

\subsection{含幺\texorpdfstring{$\Ai$}{A∞} 模}\label{sec:5p}
令 $\mathcal P=\modcat(\A^{\to})^{\mathrm{fin}}$ 为有限维、严格含幺 $\Ai$ 模组成的 $\Ai$ 范畴。有限维是指 $\M(Y_k)$ 是总维数有限的分次向量空间；严格含幺是指 $\mu^2_{\M}(b,e_{Y_k})=b$，并且在 $d\geq3$ 时，只要某个 $a_j$ 是单位元，就有 $\mu^d_{\M}(b,a_{d-1},\ldots,a_1)=0$；此外，只允许具有类似消失性质的预模同态。由于有向性，后一个条件意味着态射空间 $\hom_{\mathcal P}(\M_0,\M_1)$ 有限维。Yoneda 嵌入 $\mathcal I$ 及其在扭曲复形上的延伸 $\widetilde{\mathcal I}$，其像自然落在 $\mathcal P$ 中。另一方面，嵌入 $\mathcal P\to\Q$ 上同调全忠实；考察长度滤过所给谱序列的 $E_1$ 项即可看出：对 $\Q$ 与 $\mathcal P$，它们分别是\eref{2.14}中用于计算模 $\operatorname{Ext}$ 群的标准条形复形的非约化版本与约化版本。
\begin{corollary}\label{st:5.26}
$\mathcal P$ 是 $\A^{\to}$ 的三角化包络。
\end{corollary}
\SourcePage{95}
\begin{proof}
由上述观察，$\mathcal I:\A^{\to}\to\mathcal P$ 是上同调全忠实的 $\Ai$ 函子。此外，$\mathcal P$ 对\eref{3.32}意义下的映射锥封闭，故用与\stref{引理}{3.35}相同的论证，可知它三角化。还需证明 $\mathcal I$ 的像生成 $\mathcal P$，等价地，证明 $\widetilde{\mathcal I}:\Tw(\A^{\to})\to\mathcal P$ 是拟等价。但由\stref{引理}{5.24}，简单模位于 $\widetilde{\mathcal I}$ 的像中；另一方面，任意模都容易表示为由简单模形成的迭代锥。
\end{proof}
\begin{corollary}\label{st:5.27}
任意有向 $\Ai$ 范畴都拟同构于另一个仍然有向、而且实际上为 dg 范畴的范畴，即 $d>2$ 时 $\mu^d=0$。
\end{corollary}
\begin{proof}
$\mathcal I(Y_k)$ 在 $\mathcal P$ 中构成例外集。取相应的有向 $\Ai$ 子范畴；由于 $\mathcal P$ 严格含幺且 $\hom$ 空间有限维，这一步很直接。高阶复合映射为零是 $\Ai$ 模范畴的一般性质。
\end{proof}
当然，我们已经知道，任意 $\Ai$ 范畴都可以变换为一个拟同构的 dg 范畴（\stref{推论}{2.14}）。上述结果的要点在于：有向情形下，可以在保持 $\hom$ 空间有限维的同时做到这一点；在更一般的背景下，这是不可能的。\znote{87}{关键在于：有向 $\Ai$ 范畴的模范畴 $\mathcal P$ 既严格含幺、又只有有限维 $\hom$ 空间，故在 $\mathcal P$ 内部作形式消去（把高阶结构映射逐步杀掉）不会破坏有限维性。对一般的 $\Ai$ 范畴，把 $\mu^d$（$d>2$）消去通常要引入无穷维的 $\hom$ 空间，因此推论 5.27 的有向性假设不能去掉。[来源：并见 Seidel, Fukaya categories and Picard–Lefschetz theory, §5p（推论 5.27）与推论 2.14。]}
